%%%%%%%%%%%%%%%%%%%%%%%%%%%%%%%%%%%%%%%%
% Gautam Rajendrakumar Gare - Fixed Academic CV (2025)
% Template: AltaCV (two-column academic layout)
%%%%%%%%%%%%%%%%%%%%%%%%%%%%%%%%%%%%%%%%

\documentclass[10pt,a4paper,ragged2e,withhyper]{altacv}
\geometry{left=1cm,right=1cm,marginparwidth=8.2cm,marginparsep=1cm,top=0.8cm,bottom=0.5cm}
% \geometry{left=1cm,right=1cm,marginparwidth=8.2cm,marginparsep=1cm,top=0.8cm,bottom=0.5cm}
% \geometry{left=1cm,right=10.2cm,marginparwidth=8.2cm,marginparsep=1cm,top=0.8cm,bottom=0.5cm}

\definecolor{heading}{HTML}{0B3954}
\definecolor{accent}{HTML}{1565C0}
\definecolor{emphasis}{HTML}{283593}
\definecolor{body}{HTML}{111111}
\colorlet{name}{heading}

\hypersetup{colorlinks=true, urlcolor=accent, linkcolor=accent, citecolor=accent}
\newcommand{\publink}[2]{{\small[\href{#1}{#2}]}}

% Keep section titles and event headers on the same page as their content
\usepackage{needspace}
\NewCommandCopy{\origcvsection}{\cvsection}
\renewcommand{\cvsection}[2][]{\needspace{6\baselineskip}\origcvsection[#1]{#2}}
\NewCommandCopy{\origcvevent}{\cvevent}
\renewcommand{\cvevent}[4]{\needspace{8\baselineskip}\origcvevent{#1}{#2}{#3}{#4}}

\renewcommand{\familydefault}{\sfdefault}

\name{Gautam Rajendrakumar Gare}
\tagline{Inference-Time Adaptation \& Alignment of Vision--Language Models}
\personalinfo{
  \email{ggare@andrew.cmu.edu}
  \linkedin{gautam-gare}
  \homepage{ggare-cmu.github.io}
  \location{Pittsburgh, PA}
}

\begin{document}
\makecvheader

%%%%%%%%%%%%%%%%%%%%%%%%%%%%%%%%%%%%%%%%%%%%%%%%%%%%%%%%%%%%%%%%%%%%
% Highlighted Research Summary box
%%%%%%%%%%%%%%%%%%%%%%%%%%%%%%%%%%%%%%%%%%%%%%%%%%%%%%%%%%%%%%%%%%%%
% \vspace{-0.4em}
\vspace{0.2em}
\begin{tcolorbox}[colback=blue!5!white, colframe=accent, boxrule=0.4pt, arc=2mm, left=1mm, right=1mm, top=1mm, bottom=1mm]
% \small
\textbf{Research Summary:} 
% My research focuses on developing \textbf{causally-grounded and interpretable vision–language models (VLMs)} 
% for medical imaging. I design frameworks that align neural representations with \textbf{clinician-understandable biomarkers}, 
% using causal discovery and contrastive learning to improve reliability, transparency, and generalization across 
% patients and modalities.
% My research focuses on developing \textbf{causally-grounded, human-aligned, and interpretable 
% vision–language models (VLMs)} for object detection and medical imaging. I design frameworks that 
% align neural representations with \textbf{clinician-understandable biomarkers and human feedback}, 
% using prompt tuning, causal discovery, and contrastive learning to improve reliability, 
% transparency, and generalization across patients, domains, and modalities.
% My research focuses on developing \textbf{interpretable and human-aligned vision–language models 
% (VLMs)} for object detection and medical imaging. I design frameworks that align neural 
% representations with \textbf{clinician-understandable biomarkers and human feedback}, using 
% prompt tuning and causal discovery to improve transparency, reliability, and trustworthy 
% generalization across domains and modalities.
I make \textbf{frozen vision--language models (VLMs) do new things without touching their weights},
treating adaptation as an inference-time problem: optimizing prompts (\textbf{DetPO, ECCV 2026}),
selecting image-conditioned attributes, and reading activations (\textbf{ACL Findings 2026}).
I study when and why multimodal models generalize (\textbf{ICML 2024 Oral}) and prove these methods
in the high-stakes setting of medical imaging (5 ISBI papers, \textbf{4 U.S. patents}).
\textbf{On the job market for research scientist roles} (graduating late 2026).
\end{tcolorbox}
\vspace{0.2em}

% %%%%%%%%%%%%%%%%%%%%%%%%%%%%%%%%%%%%%%%%%%%%%%%%%%%%%%%%%%%%%%%%%%%%
% % \vspace{-0.5em}
% \cvsection{Research Summary}
% \small{My research focuses on developing interpretable, causally grounded vision–language models for medical imaging.
% I explore how contrastive and causal methods can align neural representations with human-understandable biomarkers,
% improving reliability and clinical trust.}

%%%%%%%%%%%%%%%%%%%%%%%%%%%%%%%%%%%%%%%%%%%%%%%%%%%%%%%%%%%%%%%%%%%%
\vspace{-1em}
\cvsection{Education}
\cvevent{Ph.D. in Robotics (Candidate)}{Carnegie Mellon University}{Aug 2021 -- Late 2026 (expected)}{Pittsburgh, PA}
\begin{itemize}
\item Advisors: Prof. John Galeotti, Prof. Deva Ramanan
\item Focus: Inference-time adaptation and alignment of VLMs; generalization; medical imaging AI.
\end{itemize}
\vspace{-0.5em}
\divider

\cvevent{M.S. in Electrical and Computer Engineering}{Carnegie Mellon University}{Dec 2020}{Pittsburgh, PA}
\begin{itemize}
\item GPA: 4.00/4.00 - Courses: SLAM, CV, ML, Optimization, Deep Learning
\end{itemize}
\vspace{-0.5em}

\divider

\cvevent{B.E. in Electronics and Communication}{BMS College of Engineering}{May 2016}{Bangalore, India}
% \begin{itemize}
% \item GPA: 9.4/10.0
% \end{itemize}

\vspace{-0.5em}
%%%%%%%%%%%%%%%%%%%%%%%%%%%%%%%%%%%%%%%%%%%%%%%%%%%%%%%%%%%%%%%%%%%%
\cvsection{Research Experience}

\cvevent{Soft-Prompt Tuning of Vision–Language Models for Few-Shot Object Detection}
{Robotics Institute, CMU}{Nov 2025 -- Present}{}
\begin{itemize}
    \item Designed a soft-prompting framework that learns continuous prompt embeddings 
          via gradient-based optimization, enabling efficient task adaptation without 
          modifying the underlying VLM weights.
    \item Outperformed LoRA fine-tuning in the 10-shot setting, demonstrating that 
          lightweight prompt optimization can match or exceed parameter-efficient 
          fine-tuning methods at low data regimes.
\end{itemize}

\vspace{-0.25cm}
\divider
\vspace{-0.25cm}

\cvevent{DetPO: Iterative Prompt Optimization for Few-Shot Object Detection (ECCV 2026)}{Robotics Institute, CMU}{Sep 2025 -- Jan 2026}{}
\begin{itemize}
\item Developed an iterative prompt optimization framework using false-positive and false-negative feedback to improve few-shot object detection; accepted at \textbf{ECCV 2026}.
\item Achieved state-of-the-art results with contemporary VLMs (e.g., Qwen3, Gemini3) on the RF20 object detection dataset, outperforming specialist models such as GroundingDINO under a black-box optimization setting.
\end{itemize}

\vspace{-0.25cm}
\divider
\vspace{-0.25cm}

\cvevent{Distribution-Conditioned Attribute Selection for VLMs}{Robotics Institute, CMU}{2026}{}
\begin{itemize}
\item Showed LLM-generated class descriptors are label-conditioned, not image-conditioned: removing class names collapses ImageNet accuracy from 59.5\% to 15.5\%.
\item Selecting attributes from target images in CLIP's embedding space outperforms CoOp prompt tuning with one image per class, fitting in under a minute instead of 14 hours.
\end{itemize}

\vspace{-0.25cm}
\divider
\vspace{-0.25cm}

\cvevent{Causal Feature Selection \& Natural Experiments}{Robotics Institute, CMU}{Jan 2023 -- Aug 2025}{}
\begin{itemize}
\item Designed a Markov--blanket--based causal discovery algorithm for feature selection, improving accuracy and interpretability on medical datasets.
\item Showed real-world datasets contain natural experiments exploitable via causal inference.
\end{itemize}

\vspace{-0.25cm}
\divider
\vspace{-0.25cm}

\cvevent{Generic Human–Interpretable Lung Ultrasound Biomarkers}{Biomedical Image Guidance Lab, CMU}{2021 -- 2025}{}
\begin{itemize}
\item Developed interpretable LUS biomarkers enabling causal and contrastive learning for diagnosis.
\item Showed biomarker-based features match deep features in downstream performance.
\end{itemize}

% \vspace{-0.15cm}
% \divider
% \vspace{-0.15cm}

% \cvevent{LEARNER: Contrastive Learning for Patient Progression}{CMU / ISBI 2025}{2023 -- 2025}{}
% \begin{itemize}
% \item Proposed weakly supervised contrastive pretraining for intra-patient progression learning.
% \item Improved F1 across modalities (LUS, MRI) under label-scarce settings.
% \end{itemize}

% \vspace{-0.15cm}
% \divider
% \vspace{-0.15cm}

% % \vspace{2cm}
% \cvevent{Action Recognition in Clinical Settings}{DawnLight Inc. (Stanford-affiliated)}{Jun 2021 -- Aug 2021}{}
% \begin{itemize}
% \item Built efficient RGB+IR architectures for fall detection (85\% cross-subject accuracy).
% \end{itemize}


%%%%%%%%%%%%%%%%%%%%%%%%%%%%%%%%%%%%%%%%%%%%%%%%%%%%%%%%%%%%%%%%%%%%%%%%%%%%
\cvsection{Selected Publications}

\begin{itemize}
\interlinepenalty=10000 % never split a publication entry across pages
\item \textbf{Gare, G. R.}*, N. Peri*, M. Popov, S. Jain, J. Galeotti, D. Ramanan. (*equal contribution)\\
\emph{DetPO: In-Context Learning with Multi-Modal LLMs for Few-Shot Object Detection.}\\
\textbf{European Conference on Computer Vision (ECCV)}, 2026.
\publink{https://arxiv.org/abs/2603.23455}{paper}~\publink{https://ggare-cmu.github.io/DetPO/}{project}~\publink{https://github.com/ggare-cmu/DetPO}{code}
\divider

\item \textbf{Gare, G.R.}, J. Shi, Z. Lin, D. Pathak, J. Galeotti, D. Ramanan.\\
\emph{Attributes Should Come from Images, Not Class Names: Distribution-Conditioned Attribute Selection for VLMs.}\\
\textbf{PFATCV Workshop at ECCV}, 2026. \textbf{(Oral)}
\publink{https://arxiv.org/abs/2607.18695}{paper}~\publink{https://ggare-cmu.github.io/AttributeSelect/}{project}
\divider

\item R. H. Abhinandan, J. Galeotti, D. Ramanan, \textbf{G.R. Gare}.\\
\emph{Ask Twice, Look Twice: Prompt Echoing Resolves the Question-First Paradox in Vision-Language Models.}\\
\textbf{eXCV Workshop at ECCV}, 2026.
\publink{https://arxiv.org/abs/2607.15565}{paper}~\publink{https://rakshanda-cmu.github.io/ask-twice-look-twice/}{project}
\divider

\item T. Chai, C. Mitra, B. Huang, \textbf{G.R. Gare}, Z. Lin, A. Arbelle, L. Karlinsky, R. Feris, T. Darrell, D. Ramanan, R. Herzig.\\
\emph{Activation reward models for few-shot model alignment.}\\
\textbf{Findings of the Association for Computational Linguistics (ACL)}, 2026.
\publink{https://aclanthology.org/2026.findings-acl.1709/}{paper}
\divider

\item J. Shi, \textbf{G.R. Gare}, J. Tian, S. Chai, Z. Lin, A. Vasudevan, D. Feng, F. Ferroni, S. Kong.\\
\emph{LCA-on-the-Line: Benchmarking Out-of-Distribution Generalization with Class Taxonomies.}\\
\textbf{International Conference on Machine Learning (ICML)}, 2024. \textbf{(Oral)}
\publink{https://proceedings.mlr.press/v235/shi24d.html}{paper}~\publink{https://elvishelvis.github.io/papers/lca/}{project}
\divider

\item \textbf{Gare, G.R.}, J. M. Galeotti, M. C. Mozer, D. Ramanan, N. R. Ke.\\
\emph{Do Real-World Datasets Contain Natural Experiments? An Empirical Study Using Causal Feature Selection.}\\
\textbf{arXiv:2606.03251}, 2026.
\publink{https://arxiv.org/abs/2606.03251}{paper}
\divider

\item \textbf{Gare, G.R.}, T. Fox, P. Lowery, K. Zamora, H. Tran, L. Hutchins, D. Montgomery, A. Krishnan, D. Ramanan, R. Rodriguez, B. deBoisblanc, J. Galeotti.\\
\emph{Human-Aligned Learning of Interpretable Biomarkers for Lung Ultrasound Diagnosis.}\\
\textbf{IEEE International Symposium on Biomedical Imaging (ISBI)}, 2026.
\publink{https://arxiv.org/abs/2206.08398}{paper}
\divider

\item J. Armouti, N. Madaan, R. Panda, T. Fox, L. Hutchins, A. Krishnan, R. Rodriguez, B. DeBoisblanc, D. Ramanan, J. Galeotti, \textbf{G.R. Gare}.\\
\emph{LEARNER: Learning Granular Labels from Coarse Labels using Contrastive Learning.}\\
\textbf{IEEE International Symposium on Biomedical Imaging (ISBI)}, 2026.
\publink{https://arxiv.org/abs/2411.01144}{paper}
\divider

\item M. Shivaram, \textbf{G.R. Gare}, L. Hutchins, J. Duplantis, T. Deiss, T.N. Gomes, T. Tran, K.H. Patel, K. Kirby, T.H. Fox, A. Krishnan, D. Ramanan, B. DeBoisblanc, R. Rodriguez, J. Galeotti.\\
\emph{Label uncertainty for ultrasound segmentation.}\\
\textbf{IEEE International Symposium on Biomedical Imaging (ISBI)}, 2026. \textbf{(Oral)}
\publink{https://arxiv.org/abs/2508.15635}{paper}
\divider

\item \textbf{G.R. Gare}, H.V. Tran, B.P. deBoisblanc, R.L. Rodriguez, J.M. Galeotti.
\emph{Weakly supervised contrastive learning for better severity scoring of lung ultrasound.}
\textbf{IEEE International Symposium on Biomedical Imaging (ISBI)}, 2025.
\publink{https://arxiv.org/abs/2201.07357}{paper}
\divider

\item \textbf{Gare, G.R.}, T. Fox, B. Chansangavej, A. Krishnan, R. Rodriguez, B. deBoisblanc, D. Ramanan, J. Galeotti.\\
\emph{Improving Model’s Interpretability and Reliability using Biomarkers.}\\
\textbf{Bern Interpretable AI Symposium (BIAS)}, 2023.
\publink{https://arxiv.org/abs/2402.12394}{paper}
\divider

\item \textbf{Gare, G.R.}, J. Li, R. Joshi, R. Magar, M. P. Vaze, M. Yousefpour, R. Rodriguez, J. Galeotti.\\
\emph{W-Net: Dense and diagnostic semantic segmentation of subcutaneous and breast tissue in ultrasound.}\\
\textbf{Medical Image Analysis (MIA)}, 76:102326, 2022.
\publink{https://doi.org/10.1016/j.media.2021.102326}{paper}
\divider

\item \textbf{G.R. Gare}, A. Schoenling, V. Philip, H.V. Tran, B.P. deBoisblanc, R.L. Rodriguez, J.M. Galeotti.
\emph{Dense pixel-labeling for reverse-transfer and diagnostic learning on lung ultrasound for COVID-19 and pneumonia detection.}
\textbf{IEEE International Symposium on Biomedical Imaging (ISBI)}, 2021, pp. 1406–1410.
\publink{https://ieeexplore.ieee.org/document/9433826}{paper}
\divider
\end{itemize}

%%%%%%%%%%%%%%%%%%%%%%%%%%%%%%%%%%%%%%%%%%%%%%%%%%%%%%%%%%%%%%%%%%%%%%%%%%%%
% \vspace{0.2cm}
\needspace{11\baselineskip}
\vspace{-0.5em}
\cvsection{Professional Experience}

\cvevent{AI Research Intern}{Robert Bosch LLC (Bosch Research)}{May 2026 -- Aug 2026}{Pittsburgh, PA}
\begin{itemize}
\item Researched grounding of vision--language--action (VLA) models for autonomous-vehicle applications, connecting language-conditioned action reasoning to visual scene evidence.
\end{itemize}

\divider

\cvevent{Computer Vision Research Intern}{DawnLight Inc.}{Jun 2021 -- Aug 2021}{Palo Alto, CA}
\begin{itemize}
\item Developed efficient multimodal action recognition networks (RGB + IR) for clinical monitoring.
\end{itemize}

\divider

\cvevent{Senior Data Scientist \& Software Engineer}{Sling Media R\&D}{Jul 2016 -- Jul 2019}{Bangalore, India}
\begin{itemize}
\item Led analytics pipeline for user churn modeling and AirTV Player Android App interface optimization.
\item Improved image caching, reducing startup time by 20\%.
\end{itemize}

\divider

\cvevent{Summer Research Fellow}{IIEST, Shibpur}{Jun 2015 -- Jul 2015}{India}
\begin{itemize}
\item Built real-time tele-operated robotic arm system with RGB-D feedback.
\end{itemize}


%%%%%%%%%%%%%%%%%%%%%%%%%%%%%%%%%%%%%%%%%%%%%%%%%%%%%%%%%%%%%%%%%%%%%%%%%%%%
\cvsection{Selected Patents}

{\footnotesize
\begin{itemize}
\item US12154354B2 – \textbf{System and Method for Labeling Ultrasound Data}, granted 2024.
\item US20240177000A1 – \textbf{System and Method for Training Machine-Learning Models with Probabilistic Confidence Labels}, 2024.
\item US20240177445A1 – \textbf{System and Method for Direct Diagnostic and Prognostic Semantic Segmentation of Images}, 2024.
\item US20240249498A1 – \textbf{System, Method, and Computer Program Product for Extracting Features from Imaging Biomarkers with ML Models}, 2024.
\end{itemize}
}

%%%%%%%%%%%%%%%%%%%%%%%%%%%%%%%%%%%%%%%%%%%%%%%%%%%%%%%%%%%%%%%%%%%%%%%%%%%%

\cvsection{Fellowships \& Awards}

\begin{itemize}
\item Liang Zhao Endowed Fellowship, CMU (2025)
\item CMLH Translational Fellowship, CMU (2022)
\item NSF XSEDE Compute Grant (130,500 GPU hours, 2020)
\item Neocortex EUP Compute Award, NSF PSC (2020)
\end{itemize}


%%%%%%%%%%%%%%%%%%%%%%%%%%%%%%%%%%%%%%%%%%%%%%%%%%%%%%%%%%%%%%%%%%%%%%%%%%%%
% \clearpage
%%%%%%%%%%%%%%%%%%%%%%%%%%%%%%%%%%%%%%%%%%%%%%%%%%%%%%%%%%%%%%%%%%%%%%%%%%%%
% RIGHT COLUMN


%%%%%%%%%%%%%%%%%%%%%%%%%%%%%%%%%%%%%%%%%%%%%%%%%%%%%%%%%%%%%%%%%%%%%%%%%%%%
\cvsection{Teaching \& Service}

\begin{itemize}
\item Guest Lecturer – 16-720 Computer Vision (Fall 2026)
\item Research Mentor – supervised the prompt-echoing VLM study (arXiv:2607.15565, 2026)
\item Research Mentor – supervised the lung-ultrasound CHF-readmission study (arXiv:2605.18878, 2026)
\item Teaching Assistant – 16-720 Computer Vision (Spring 2022, Fall 2021)
\item Reviewer – ECCV (2026), eXCV Workshop at ECCV (2026), MICCAI (2022–2024), HSMR (2023)
\item Invited Panelist – LLCOVID19 MICCAI Workshop (2021)
\item Mentorship – CMU AI Mentoring Program (2021)
\end{itemize}

%%%%%%%%%%%%%%%%%%%%%%%%%%%%%%%%%%%%%%%%%%%%%%%%%%%%%%%%%%%%%%%%%%%%%%%%%%%%
\cvsection{Media Coverage}

\begin{itemize}
\item Megan Harris, \emph{“Seeing Beneath The Surface”}, \textbf{LINK Magazine},
School of Computer Science, Carnegie Mellon University, Summer 2022 Issue.
[\href{https://magazine.cs.cmu.edu/seeing-beneath-the-surface}{article}]
\item \emph{“CMU and others propose a highly practical `model generalization' metric, improving out-of-distribution accuracy by 6\% | ICML 2024 Oral”}, \textbf{New Zhiyuan}, December 26, 2024, covering \emph{LCA-on-the-Line}.
[\href{https://mp.weixin.qq.com/s/MC8eDigX5WOJySh8StkBDA}{article}]
\end{itemize}

\end{document}